%!TEX TS-program = xelatex
\documentclass[]{friggeri-cv}
\usepackage{afterpage}
\usepackage{hyperref}
\usepackage{color}
\usepackage{xcolor}
\usepackage{smartdiagram}
\usepackage{fontspec}
% if you want to add fontawesome package
% you need to compile the tex file with LuaLaTeX
% References:
%   http://texdoc.net/texmf-dist/doc/latex/fontawesome/fontawesome.pdf
%   https://www.ctan.org/tex-archive/fonts/fontawesome?lang=en
%\usepackage{fontawesome}
\usepackage{metalogo}
\usepackage{dtklogos}
\usepackage[utf8]{inputenc}
\usepackage{tikz}
\usetikzlibrary{mindmap,shadows}
\hypersetup{
pdftitle={Alex Irmel Oviedo Solis -  Computer Engineer},
pdfauthor={Ing. Alex Irmel Oviedo Solis},
pdfsubject={Curriculum vitae},
pdfkeywords={},
colorlinks=false,           % no lik border color
allbordercolors=white       % white border color for all
}
\smartdiagramset{
bubble center node font = \footnotesize,
bubble node font = \footnotesize,
% specifies the minimum size of the bubble center node
bubble center node size = 0.5cm,
%  specifies the minimum size of the bubbles
bubble node size = 0.5cm,
% specifies which is the distance among the bubble center node and the other bubbles
distance center/other bubbles = 0.3cm,
% sets the distance from the text to the border of the bubble center node
distance text center bubble = 0.5cm,
% set center bubble color
bubble center node color = pblue,
% define the list of colors usable in the diagram
set color list = {lightgray, materialcyan, orange, green, materialorange, materialteal, materialamber, materialindigo, materialgreen, materiallime},
% sets the opacity at which the bubbles are shown
bubble fill opacity = 0.6,
% sets the opacity at which the bubble text is shown
bubble text opacity = 0.5,
}

\addbibresource{bibliography.bib}
\RequirePackage{xcolor}
\definecolor{pblue}{HTML}{0395DE}

\begin{document}
\header{Alex Irmel}{Oviedo Solis}
{Computer Engineer}

% Fake text to add separator
\fcolorbox{white}{gray}{\parbox{\dimexpr\textwidth-2\fboxsep-2\fboxrule}{%
.....
}}

% In the aside, each new line forces a line break
\begin{aside}
    \includegraphics[scale=0.18]{img/aoviedo.png}
    \section{Direcci\'on}
    Av. de la cultura 2505
    Cusco, Cusco
    ~
    \section{Telefonos}
    +51 930 328 402
    ~
    \section{Email}
    \href{mailto:aoviedosolis@gmail.com}{\textbf{aoviedosolis@}\\gmail.com}
    \href{mailto:alexove@fedoraproject.org}{\textbf{alexove@}\\fedoraproject.org}
    ~
    \section{Web \& Git}
    \href{https://github.com/alexove}{https://github.com/alexove}
    % ~
    % % use  \hspace{} or \vspace{} to change bubble size, if needed
    % \section{Cualidades Personales}
    % \smartdiagram[bubble diagram]{
    % \textbf{Trabajo}\\\textbf{en equipo},
    % \textbf{Iniciativa},
    % \textbf{Curiosidad},
    % \textbf{Resolucion}\\\textbf{problemas},
    % \textbf{\vspace{2mm}Gerencia\vspace{2mm}},
    % \textbf{Organizado}
    % }
    ~
    \section{Languages}
    \textbf{Spanish}\includegraphics[scale=0.40]{img/5stars.png}
    \textbf{English}\includegraphics[scale=0.40]{img/3stars.png}
    \textbf{Russian}\includegraphics[scale=0.40]{img/2stars.png}
    ~
\end{aside}
~
\section{Computer engineer}
\emph{With a predisposition to meet goals and objectives, a proactive attitude, 
the ability to work in multidisciplinary teams, teamwork under pressure, the 
ability to adapt to changes, high reliability, aspiration for professional 
development and the ability to interact with people inside and outside the 
organization and a desire for continuous progress.}
\\
\section{Education}
\begin{entrylist}
    \entry
    {2019 - 2021}
    {Master in Information Security}
    {Universidad de la Rioja - Sede Mexico}
    {Master's degree focused on the latest protection techniques against vulnerabilities 
     in operating systems, software, databases, web systems and all the risks inherent 
     to the use of Information Technologies and their possible impact within the public 
     and private organizational structure..\\}
    \entry
    {2016 - 2017}
    {MBA-not graduated}
    {Universidad Andina del Cusco}
    {The Master's in Business Administration is a program aimed at training professionals 
	from the Cusco Region and the country, engaging them with the reality of an environment 
	that increasingly demands knowledge and the correct and adequate use of information, 
	with the ability to respond with rigor, opportunity and social relevance to local, 
	regional and national problems.\\}
    \entry
    {2009 - 2012}
    {Computer Engineering Degree}
    {Universidad Nacional de San Antonio Abad del Cusco}
    {Honest, ethical, caring professional with a solid background in basic sciences, IT and 
	systems. Competent in the development of IT projects, with the ability to design, 
	manage, audit and monitor their progress.\\}
    \entry
    {2002 - 2009}
    {Bachelor in Computer Engineering}
    {Universidad Nacional de San Antonio Abad del Cusco}
    {Honest, ethical, caring professional with a solid background in basic sciences, IT and 
	systems. Competent in the development of IT projects, with the ability to design, 
	manage, audit and monitor their progress.}
\end{entrylist}
\\

\newpage
\begin{aside}
    ~
    ~
    \section{Programming Languages}
    \includegraphics[scale=0.20]{img/java}
    \includegraphics[scale=0.07]{img/groovy}
    \includegraphics[scale=0.3]{img/bash}
    ~
    ~
    \section{Operating Systems}
    \textbf{GNU/Linux}\includegraphics[scale=0.40]{img/5stars.png}
    \textbf{Windows}\includegraphics[scale=0.40]{img/1stars.png}
    ~
    \section{Data bases}
    ~
    \includegraphics[scale=0.04]{img/mysql.png}
    \includegraphics[scale=0.30]{img/mariadb.png}
    ~
    \section{Web development}
    ~
    \includegraphics[scale=0.02]{img/grails.png}
    ~
\end{aside}
\section{Certifications}
\begin{entrylist}
    \entry
    {04/2016}
    {Linux Foundation Certified System Administrator LFCS}
    {Linux Foundation}
    {Linux Foundation certifications give you a way to differentiate yourself in
    a job market that's hungry for your skills. We've taken a new, innovative
    approach to Linux certification that allows you to showcase your skills in a
     way that other sysadmins will respect and employers will trust.\\
     \emph{Certificate ID N°{:} LFCS-1600-0744-0100}}
\end{entrylist}

\section{Work experience}
\begin{entrylist}
    \entry
    {03/24 - Ahora}
    {Data Center Manager}
    {Cusco Regional Goverment}
	{Performed activities: GNU/Linux servers administration(CentOS, Debian, Ubuntu), network monitoring,
	SIAF-MEF support and maintenance, Digital Goverment and Digital Security (early stage).}
    \entry
    {09/23 - 02/24}
    {BIM developer}
    {Consorcio Rios del Norte S.A.C.}
    {Performed activities: Planning, development of quality management systems and process automation within the BIM metodology.}
    \entry
    {12/22 - 08/23}
    {Head of IT and Systems}
    {MacSalud S.A.C.}
    {Performed activities: Planning, develpment, resource and task allocation, coordination with information systems providers
     network monitoring and support, data center management and Internet.}
    \entry
    {02/22 - 11/22}
    {Coordinador de proyecto de software}
    {Regional Education Management of Cusco}
    {Performed activities: Planning, development, resource and tasks allocation,
     implementation, monitoring and delivery of Virtual Education Platform.}
    \entry
    {08/21 - 09/21}
    {Information and Apps Security Analist}
    {Caja Municipal Cusco}
    {Performed activities: Information security controls monitoring. information security risk analisis.}
    \entry
    {03/21 - 07/21}
    {Data Center Manager}
    {Cusco Regional Goverment}
	{Performed activities: GNU/Linux servers management (CentOS, Debian, Ubuntu), network monitoring and support, SIAF-MEF 
	and SIGA-MEF support and maintenance.\\}
    \entry
    {01/20 - 03/21}
    {Head of IT and Systems}
    {Cusco Regional Goverment}
    {Performed activities: Software project coordination, network monitoring and support, data center monitoring and support,
	Digital Goverment.\\}
    \entry
    {01/18 - 12/19}
    {Data center Manager}
    {Cusco Regional Goverment}
    {Performed activities: GNU/Linux servers management (CentOS, Debian, Ubuntu), network monitoring and support, SIAF-MEF 
	and SIGA-MEF support and maintenance.\\}
    \entry
    {02/17 - 06/17}
    {GIS developer}
    {CEC Guaman Poma de Ayala}
    {GIS developer with java, gvSIG and PostgreSQL.\\}
\end{entrylist}
\begin{entrylist}
    \entry
    {04/16 - 07/16}
    {Data center Manager}
    {Cusco Regional Goverment}
    {Performed activities: GNU/Linux servers management (CentOS, Debian, Ubuntu), network monitoring and support, SIAF-MEF 
	and SIGA-MEF support and maintenance.\\}
    \entry
    {02/16 - 03/16}
    {Microsoft Office Teacher}
    {CENFOTI – Universidad Andina del Cusco.}
    {Performed activities: Microsoft Office Teacher.\\}
    \entry
    {07/15 - 12/15}
    {Data center Manager}
    {Cusco Regional Goverment}
    {Performed activities: GNU/Linux servers management (CentOS, Debian, Ubuntu), network monitoring and support, SIAF-MEF 
	and SIGA-MEF support and maintenance.\\}
    \entry
    {01/15 - 03/15}
    {Data center Manager}
    {Cusco Regional Goverment}
    {Performed activities: GNU/Linux servers management (CentOS, Debian, Ubuntu), network monitoring and support, SIAF-MEF 
	and SIGA-MEF support and maintenance.\\}
    \entry
    {06/14 - 10/14}
    {Software Developer}
    {Las Bambas - ManPower}
	{Software Developer for cooper mining operations (drilling and blasting).\\}
    \entry
    {07/12 - 02/14}
    {Java web developer}
    {Java For Smart Information Technologies - Cusco}
    {Java web developer, JSP and basic GWT, Android development, GRAILS Framework, VPS management.\\}
    \entry
    {11/12 - 01/13}
    {Data Center associate}
    {{RootWay Internet Services \& Consulting - Uruguay}}
    {{Administraci\'on y migraci\'on de servidores virtuales, GNU/Linux en AmazonWS. \\ }}
    \entry
    {01/12 - 05/12}
    {Jefe de la Oficina de Inform\'atica.}
    {Municipalidad Provincial de Quispicanchi}
    {Jefe de la Oficina de Inform\'atica.\\}
\end{entrylist}
\begin{aside}
    ~
    \section{GNU/Linux}
    ~
    \includegraphics[scale=0.12]{img/fedora}
    \includegraphics[scale=0.12]{img/centos}
    ~
\end{aside}
\section {Diplomados}
\begin{entrylist}
  \entry
  {10/2020 - 03/2021}
  {Gestión de la Seguridad y Marco Legal}
  {Universidad Internacional de la Riona en México}
  {Realizado del 26 de octubre de 2020 al 7 de marzo de 2021, con una duración de 432 horas}
  \entry
  {05/2020 - 08/2020}
  {Seguridad en los Nuevos Entornos y Auditoría}
  {Universidad Internacional de la Riona en México}
  {Realizado del 4 de mayo de 2020 al 30 de agosto de 2020, con una duración de 360 horas}
  \entry
  {10/2019 - 03/2020}
  {Seguridad en Redes, Sistemas y Aplicaciones}
  {Universidad Internacional de la Riona en México}
  {Realizado del 28 de octubre de 2019 al 8 de marzo de 2020, con una duración de 360 horas}
\end{entrylist}
\newpage
\section{Cursos y seminarios}
\begin{entrylist}
    \entry
    {08/23 - 08/23}
    {Habilidades Gerenciales para el Pequeño Empresario}
    {UDEMY}
    {Curso dictado por Tomas Bradanovic durante 2 horas lectivas.\\}
    \entry
    {01/23 - 01/23}
    {Asterisk nativo, configura tu Central Telefonica desde cero}
    {UDEMY}
    {Curso dictado por Roger Mamani durante 11 horas lectivas.\\}
    \entry
    {01/23 - 01/23}
    {Gestión de riesgos ISO/IEC 27005, con formatos editables}
    {UDEMY}
    {Curso dictado por Israel Rogelio Gomez durante 2 horas lectivas.\\}
    \entry
    {01/23 - 01/23}
    {ISO/IEC 27001 Implementando Seguridad de la Información}
    {UDEMY}
    {Curso dictado por Grover Vasquez durante 2 horas lectivas.\\}
    \entry
    {11/21 - 11/21}
    {Diseño y construcción de redes de fibra optica}
    {UDEMY}
    {Curso ofrecido por Gaston Fenske durante 2.5 horas lectivas.\\}
    \entry
    {08/21 - 08/21}
    {Analisis de Vulnerabilidades Digitales}
    {Centro Nacional de Seguridad Digital}
    {Curso ofrecido por el Centro Nacional de Seguridad Digital (PCM-SEGDI) durante 24 horas lectivas.\\}
    \entry
    {06/21 - 06/21}
    {Auditor Interno ISO 27001}
    {Centro Nacional de Seguridad Digital}
    {Curso ofrecido por el Centro Nacional de Seguridad Digital (PCM-SEGDI) durante 24 horas lectivas.\\}
    \entry
    {05/21 - 05/21}
    {Linux Essentials}
    {Centro Nacional de Seguridad Digital}
    {Curso ofrecido por el Centro Nacional de Seguridad Digital (PCM-SEGDI) durante 24 horas lectivas.\\}
    \entry
    {04/21 - 04/21}
    {SCRUM Práctico en Proyectos de Software}
    {UDEMY}
    {Curso dictado por Hector Bravo durante 4.5 horas lectivas.\\}
    \entry
    {10/18 - 10/18}
    {Curso PRE-MTCNA Mikrotik}
    {COMUTEL}
    {Curso dictado por COMUTEL durante 16 horas lectivas.\\}
    \entry
    {12/15 - 01/16}
    {Curso OFICIAL ITIL FOUNDATION}
    {BSGrupo}
    {Realizado entre el 06 de diciembre del 2015 y el 06 de enero del 2016.\\}
\end{entrylist}

\section{Congresos y ponencias}
\begin{entrylist}
  \entry
    {11/2016}
    {\textbf{Ponente} VI Encuentro de Ing. de Sistemas e Inform\'atica }
    {UTEA}
    {Realizado del 16 al 18 de noviembre del 2015 en la Universidad Tecnológica de los Andes - Sede Cusco.\\}
\end{entrylist}
\begin{entrylist}
    \entry
    {09/2015}
    {\textbf{Ponente} Fedora Users and Developers Conference - Argentina}
    {Fedora Community}
    {Realizando en la ciudad de Cordoba-Argentina, entre 10 y 12 de septiembre
    del 2015 en la Universidad Nacional de C\'ordoba – Argentina.\\}
    \entry
    {08/2015}
    {\textbf{Ponente} {XXIII CONEISC}}
    {Universidad Catolica Santa Maria}
    {Realizado entre el 17 al 21 de agosto del 2015 en la Universidad Catolica Santa Maria, Arequipa, Peru.\\}
    \entry
    {11/2014}
    {\textbf{Ponente} VI Festival Internacional de Software Libre SOFLISEOANE}
    {Instituto de Educaci\'on Superior Tecnol\'ogico Publico Manuel Seoane Corrales}
    {Realizado el 28 y 29 de noviembre del 2014 en el Instituto de Educaci\'on Superior
    Tecnol\'ogico Publico Manuel Seoane Corrales, Lima.\\}
    \entry
    {04/2014}
    {\textbf{Ponente} 15-vo Foro Internacional do Software Livre-Porto Alegre-Brasil}
    {FISL}
    {Realizado entre el 7 y 10 de mayo del 2014 en la ciudad de Porto Alegre-Brasil.\\}
    \entry
    {01/14 - 03/14}
    {Curso Taller de "SEGURIDAD INFORM\'ATICA"}
    {Grupo Development In Peru SCRL}
    {Organizado por el Grupo Development In Peru SCRL.\\}
    \entry
    {12/2015}
    {\textbf{Ponente} Semana Informatica INSANE CODE}
    {IST Tupac Amaru}
    {Organizado por el Instituto de Educaci\'on Superior Tecnol\'ogico Publico "Tupac
    Amaru" del Cusco.}
\end{entrylist}
\begin{entrylist}
    \entry
    {10-2013}
    {\textbf{Ponente} IV ENCUENTRO DE INGENIERIA DE SISTEMAS E INFORMATICA}
    {UTEA}
    {Realizado los d\'ias 26 y 27 de octubre del 2013 organizado por la Universidad
    Tecnol\'ogica de los Andes.\\}
    \entry
    {01/13 - 03/13}
    {Curso ANDROID{:} PROGRAMACI\'ON DE APLICACIONES}
    {UPV}
    {Realizado del 10 de enero del 2013 hasta el 30 de marzo del 2013, dado por
    la Universidad Polit\'ecnica de Valencia-Espa\~na.}
    \entry
    {09-2012}
    {Encuentro GOOGLE DEVELOPMENT FEST LIMA}
    {Google Developers Group}
    {Realizado el 27 y 28 de octubre del 2012 y organizado por Google Developer Group Lima.},
    \entry
    {08-2012}
    {Seminario Taller de Regional ESTRATEGIA NACIONAL DE GOBIERNO ELECTRONICO}
    {ONGEI}
    {Realizado el 13 y 14 de septiembre del 2012, organizado por la Oficina
    Nacional de Gobierno Electr\'onico e Inform\'atica.}
\end{entrylist}
\section{Otras actividades}
\begin{entrylist}
    \entry
    {08-2013}
    {Organizador de FUDCon Cusco 2013}
    {Fedora Communitty, UNSAAC}
    {Fedora Users and Developers Conference Latin América Cusco 2013, organizado
    por la comunidad Fedora Perú, Universidad Nacional de San Antonio Abad del
    Cusco y Municipalidad Provincial del Cusco, realizado el 26,27 y 28 de
    septiembre del 2013.}
    \entry
    {03/13 - ahora}
    {Embajador Fedora}
    {Fedora Community}
    {Ambassadors are the representatives of Fedora. Ambassadors ensure the public
    understand Fedora's principles and the work that Fedora is doing. Additionally
    Ambassadors are responsible for helping to grow the contributor base, and to
    act as a liaison between other FLOSS projects and the Fedora community. }
\end{entrylist}
%
%\begin{flushleft}
%    \emph{May 8th, 2016}
%\end{flushleft}
%\begin{flushright}
%    \emph{John Snow}
%\end{flushright}

\end{document}
